\documentclass{BeamerTemplate}

\title{Module 2 - Probabilités}

\subtitle{STT-1920 Méthodes statistiques}
\date{\today}

\begin{document}
\begin{frame}[t]
  \titlepage
\end{frame}

%%%%%%%%%%%%%%%%%%%%%%%%%%%%%%%%%%%
\section{Expérience aléatoire}
%%%%%%%%%%%%%%%%%%%%%%%%%%%%%%%%%%%

\begin{frame}[t]{Expérience aléatoire}

\begin{Boite}{Expérience aléatoire $\mathcal{E}$}
Expérience dont le résultat dépend du hasard et ne peut être prédit avec certitude.
\end{Boite}

\medskip

\begin{Boite}{Espace échantillonnal $\mathcal{U}$}
Ensemble de tous les résultats possibles de $\mathcal{E}$.
\end{Boite}

\medskip
\underline{Exemples :}
\begin{itemize}\itemsep3mm
  \item Lancer un dé à 6 faces et noter le résultat : $\mathcal{U} = \{1, 2, 3, 4, 5, 6\}$.
  \item Mesurer le diamètre d'une pièce usinée : $\mathcal{U} = (0, +\infty)$.
  \item Compter le nombre de pannes en une semaine : $\mathcal{U} = \{0, 1, 2, \ldots\}$.
\end{itemize}

\end{frame}

%%%%%%%%%%%%%%%%%%%%%%%%%%%%%%%%%%%
\section{Événements}
%%%%%%%%%%%%%%%%%%%%%%%%%%%%%%%%%%%

\begin{frame}[t]{Événements}

\begin{Boite}{Événement}
Sous-ensemble de l'espace échantillonnal $\mathcal{U}$, soit un ensemble de résultats possibles.
\end{Boite}

\medskip
Opérations sur les événements $A$ et $B$ :
\begin{itemize}\itemsep5mm
  \item \textbf{Complémentaire} $\overline{A}$ : tous les résultats qui ne sont \emph{pas} dans $A$.
  \item \textbf{Intersection} $A \cap B$ : résultats qui sont à la fois dans $A$ \emph{et} dans $B$.
  \item \textbf{Union} $A \cup B$ : résultats qui sont dans $A$ \emph{ou} dans $B$ (ou les deux).
\end{itemize}

\end{frame}

\begin{frame}[t]{Événements mutuellement exclusifs}

\begin{Boite}{Événements mutuellement exclusifs}
$A$ et $B$ sont mutuellement exclusifs (ou incompatibles) si $A \cap B = \emptyset$.
\end{Boite}

\medskip
Cela signifie que $A$ et $B$ \alert{ne peuvent pas se réaliser simultanément}.

\bigskip
\underline{Exemple :} On pige une carte dans un jeu standard.
\begin{itemize}\itemsep3mm
  \item $A$ = «la carte est un cœur» et $B$ = «la carte est un trèfle».\\
  $\Rightarrow$ $A$ et $B$ sont mutuellement exclusifs.
  \item $A$ = «la carte est un cœur» et $C$ = «la carte est un roi».\\
  $\Rightarrow$ $A$ et $C$ ne sont pas mutuellement exclusifs (roi de cœur).
\end{itemize}

\end{frame}

%%%%%%%%%%%%%%%%%%%%%%%%%%%%%%%%%%%
\section{Probabilité}
%%%%%%%%%%%%%%%%%%%%%%%%%%%%%%%%%%%

\begin{frame}[t]{Définition de la probabilité}

\begin{Boite}{Probabilité d'un événement}
$P(A)$ est un nombre dans $[0, 1]$ mesurant la \alert{vraisemblance} que $A$ se réalise.
\end{Boite}

\medskip
Trois axiomes fondamentaux (Kolmogorov) :
\begin{enumerate}\itemsep4mm
  \item $0 \le P(A) \le 1$ pour tout événement $A$.
  \item $P(\mathcal{U}) = 1$.
  \item Si $A$ et $B$ sont mutuellement exclusifs : $P(A \cup B) = P(A) + P(B)$.
\end{enumerate}

\end{frame}

\begin{frame}[t]{Règles de calcul}

\begin{Boite}{Complément}
$$P(\overline{A}) = 1 - P(A)$$
\end{Boite}

\begin{Boite}{Addition générale}
$$P(A \cup B) = P(A) + P(B) - P(A \cap B)$$
\end{Boite}

\medskip
\underline{Note :} Si $A$ et $B$ sont mutuellement exclusifs ($P(A \cap B) = 0$) :
$$P(A \cup B) = P(A) + P(B)$$

\end{frame}

\begin{frame}[t]{Exemple — groupes sanguins}

Les proportions dans la population québécoise sont :

\begin{center}
\begin{tabular}{|c|c|c|c|c|}
\hline
Groupe & A & B & AB & O \\ \hline
Proportion & 0{,}43 & 0{,}08 & 0{,}03 & 0{,}46 \\ \hline
\end{tabular}
\end{center}

\bigskip
Un individu est choisi au hasard. Quelle est la probabilité qu'il appartienne au groupe B ou au groupe O ?

\uncover<2->{
\bigskip
Les groupes sanguins sont mutuellement exclusifs :
$$P(B \cup O) = P(B) + P(O) = 0{,}08 + 0{,}46 = \alert{0{,}54}$$
}

\end{frame}

%%%%%%%%%%%%%%%%%%%%%%%%%%%%%%%%%%%
\section{Probabilité conditionnelle et indépendance}
%%%%%%%%%%%%%%%%%%%%%%%%%%%%%%%%%%%

\begin{frame}[t]{Probabilité conditionnelle}

\begin{Boite}{Probabilité conditionnelle}
La probabilité de $A$ sachant que $B$ s'est réalisé :
$$P(A \mid B) = \frac{P(A \cap B)}{P(B)}, \quad P(B) > 0$$
\end{Boite}

\medskip
La connaissance de $B$ \alert{modifie} la probabilité de $A$.

\bigskip
\begin{Boite}{Règle de multiplication}
$$P(A \cap B) = P(A \mid B)\,P(B) = P(B \mid A)\,P(A)$$
\end{Boite}

\end{frame}

\begin{frame}[t]{Événements indépendants}

\begin{Boite}{Indépendance}
$A$ et $B$ sont indépendants si la réalisation de l'un n'affecte pas la probabilité de l'autre :
$$P(A \mid B) = P(A) \quad \Longleftrightarrow \quad P(A \cap B) = P(A)\,P(B)$$
\end{Boite}

\medskip
\underline{Attention :} indépendance $\neq$ exclusivité mutuelle.
\begin{itemize}\itemsep3mm
  \item Deux événements exclusifs ne peuvent pas se réaliser ensemble $\Rightarrow$ la réalisation de l'un \emph{exclut} l'autre : ce n'est pas de l'indépendance (à moins que $P(A)=0$ ou $P(B)=0$).
\end{itemize}

\end{frame}

%%%%%%%%%%%%%%%%%%%%%%%%%%%%%%%%%%%
\section{Loi des probabilités totales}
%%%%%%%%%%%%%%%%%%%%%%%%%%%%%%%%%%%

\begin{frame}[t]{Loi des probabilités totales}

Soit $B_1, B_2, \ldots, B_k$ une \alert{partition} de $\mathcal{U}$ :
$$\mathcal{U} = B_1 \cup B_2 \cup \cdots \cup B_k, \quad B_i \cap B_j = \emptyset \text{ si } i \neq j$$

\begin{Boite}{Loi des probabilités totales}
$$P(A) = \sum_{j=1}^{k} P(A \mid B_j)\,P(B_j)$$
\end{Boite}

\medskip
Utile quand on connaît $P(A \mid B_j)$ mais pas $P(A)$ directement.

\end{frame}

\begin{frame}[t]{Théorème de Bayes}

\begin{Boite}{Théorème de Bayes}
Avec la même partition $B_1, \ldots, B_k$ :
$$P(B_i \mid A) = \frac{P(A \mid B_i)\,P(B_i)}{\displaystyle\sum_{j=1}^{k} P(A \mid B_j)\,P(B_j)}$$
\end{Boite}

\medskip
Permet de \alert{mettre à jour} une probabilité à la lumière d'une nouvelle observation.

\bigskip
\underline{Exemple :} Test de dépistage — quelle est la probabilité d'être réellement malade si le test est positif?

\end{frame}

\end{document}
